% SPDX-License-Identifier: Apache-2.0
%
% A reservation station in TxHDL: the tutorial on the part. Built by
% //docs:station. Every listing is included from the tree and every
% printout and figure is what the build made by running the example.
\documentclass[journal,9pt]{IEEEtran}

\newcommand{\listingsize}{\fontsize{6}{7}\selectfont}
% SPDX-License-Identifier: Apache-2.0
%
% The house preamble, shared by every document in this directory. Loaded
% after \documentclass, before \title. Keeping it in one file is what
% stops two articles drifting apart in font, listing style or figure
% style.
% Computer Modern. IEEEtran selects Times (ptm) by default, so the face has
% to be chosen explicitly.
%
% Latin Modern is the Computer Modern variety used here, and the choice is
% forced rather than aesthetic. Under T1 encoding, plain `cmr` has no Type 1
% outlines unless cm-super is installed, and it is not in the pinned
% distribution, so pdflatex falls back to EC bitmaps and every glyph in the
% PDF becomes a Type 3 font. That was measured: the first build produced a
% document with 10 Type 3 fonts and no Type 1. Latin Modern is Computer
% Modern redrawn with a full T1 repertoire and real .pfb outlines, so the
% same design comes out scalable.
\usepackage[T1]{fontenc}
\usepackage[utf8]{inputenc}
\usepackage{lmodern}
% microtype is not in the pinned TeX distribution. emergencystretch alone
% keeps the two-column measure from producing overfull lines.
\setlength{\emergencystretch}{3em}

\usepackage{amsmath}
\usepackage{amssymb}
\usepackage{amsthm}
\usepackage{booktabs}
\usepackage{array}
% enumitem is not in the pinned TeX distribution; the list spacing below
% does what \setlist would have done.
\usepackage{float}
\usepackage{xcolor}
\usepackage{listings}
\usepackage{tikz}
\usetikzlibrary{arrows.meta,positioning,fit,backgrounds,calc}
\usepackage[hidelinks,bookmarks=true]{hyperref}

% Numbered, definition-style box for a rule the reader is meant to apply.
\theoremstyle{definition}
\newtheorem{rulebox}{Rule}
\newtheorem{example}{Example}

% Unnumbered asides.
\theoremstyle{remark}
\newtheorem*{tip}{Tip}
\newtheorem*{pitfall}{Pitfall}

% Tighter lists, in place of enumitem's \setlist.
\setlength{\topsep}{2pt}
\setlength{\itemsep}{1pt}
\setlength{\parsep}{0pt}

\newcommand{\code}[1]{\texttt{#1}}
\newcommand{\term}[1]{\emph{#1}}

% Syntax highlighting. Four roles, distinguishable in colour and still
% distinguishable in greyscale, because an IEEE article gets printed: the
% keyword is the only bold role, the comment is the only italic one, and the
% two remaining roles differ in lightness as well as in hue.
\definecolor{kwblue}{RGB}{20,60,150}
\definecolor{cmtgreen}{RGB}{90,120,90}
\definecolor{strred}{RGB}{150,50,40}
\definecolor{numgray}{gray}{0.45}
\definecolor{framegray}{gray}{0.70}
\definecolor{bgshade}{gray}{0.97}
% A darker fill for figure cells. The listing background has to stay very
% light to keep code readable; a figure cell has to be clearly shaded or the
% distinction it draws is invisible in print.
\definecolor{cellshade}{gray}{0.90}

% The listing size is the document's choice, because it depends on the
% column: a two-column article sets `\listingsize` to 6pt and keeps its
% listings within 70 columns, a one-column document keeps the body size
% and includes source files at 80. Lines are never broken; a line that does not fit
% overflows the frame, which is visible, rather than wrapping, which is
% not.
\providecommand{\listingsize}{\scriptsize}

% `'` is deliberately not a string delimiter: the language uses it for loop
% labels, as Rust does, so treating it as a quote would swallow the rest of
% the line.
\lstdefinestyle{txhdl}{
  basicstyle=\ttfamily\listingsize,
  keywordstyle=\bfseries\color{kwblue},
  commentstyle=\itshape\color{cmtgreen},
  stringstyle=\color{strred},
  numberstyle=\tiny\color{numgray},
  morekeywords={mod,use,pub,const,type,struct,enum,trait,impl,fn,async,let,mut,
    match,if,else,loop,while,for,in,break,continue,return,as,await,
    module,bus,channel,transaction,protocol,pipeline,flow,seq,config,
    port,stage,pipe,instance,connect,bind,tag,domain,blackbox,foreign,
    property,role,signal,handler,true,false,
    sequence,interface,modport,implements,package,import,var,wire,func,proc,
    case,default,next,fork,branch,join,state,fsm},
  morecomment=[l]{//},
  morestring=[b]",
  showstringspaces=false,
  columns=fullflexible,
  keepspaces=true,
  breaklines=false,
  % Framed, so a listing is bounded rather than floating in the column.
  frame=single,
  rulecolor=\color{framegray},
  framesep=3pt,
  backgroundcolor=\color{bgshade},
  xleftmargin=3pt,
  xrightmargin=3pt,
  aboveskip=6pt,
  belowskip=6pt,
  % Every listing is numbered and captioned, so it can be referred to by
  % number. `caption` without `float` numbers the listing and leaves it
  % where it was written, which is what a listing in running prose needs.
  captionpos=b,
  abovecaptionskip=2pt,
  belowcaptionskip=6pt,
}

% Figures drawn as text. Framed the same way, and with the highlighting
% switched off, because a box-and-arrow diagram is not code and half its
% words turning blue reads as an accident. Setting a style adds keys rather
% than clearing the ones already set, so the three roles are neutralised by
% hand rather than by leaving them out.
\lstdefinestyle{ascii}{
  basicstyle=\ttfamily\listingsize,
  keywordstyle=\ttfamily,
  commentstyle=\ttfamily,
  stringstyle=\ttfamily,
  columns=fullflexible,
  keepspaces=true,
  frame=single,
  rulecolor=\color{framegray},
  framesep=3pt,
  backgroundcolor=\color{bgshade},
  xleftmargin=3pt,
  xrightmargin=3pt,
  aboveskip=6pt,
  belowskip=6pt,
}
\lstset{style=txhdl}

% House TikZ styles. `shorten` lives on the arrow styles rather than on each
% arrow, so every arrow in the document gets the gap, including ones added
% later. TikZ clips a path at a node's border, so without it an arrowhead
% butts against the box with no gap at all. Keep `node distance` well above
% twice the shorten, or the arrow shrinks to nothing.
\tikzset{
  box/.style={draw,rounded corners=2pt,align=center,inner sep=4pt,
              font=\footnotesize},
  boxf/.style={box,fill=bgshade},
  lbl/.style={font=\scriptsize\itshape,align=center},
  fl/.style={-{Stealth[length=4pt]},shorten >=2.5pt,shorten <=2.5pt},
  fld/.style={fl,dashed},
}



\InputIfFileExists{buildstamp.tex}{}{}
\providecommand{\buildstamp}{Draft build (outside the Bazel flow).}

\title{A Reservation Station in TxHDL}

\author{Filip Filmar and Dragiša Janković%
\thanks{Both authors are independent. Filip Filmar is the
corresponding author, at f@hdlfactory.com; Dragiša Janković is
reached at linkedin.com/in/dragisa-jankovic.}%
\thanks{This is the tutorial on one part of the language's library:
the reservation station, for a reader who has met TxHDL through the
step-by-step tutorial~\cite{tutorial} and wants to use the part, or
to read how it is made. Every listing is included from the project
repository \code{HDL/txhdl} at build time and every printout and
figure is what the build produced by running the example. The part's
source is in the runtime document~\cite{runtime}, the example in the
examples document~\cite{examples}, and the cover~\cite{cover} lists
every document.}%
\thanks{The authors used a large language model, Claude, as an
assistant in exploring the concepts and in writing and constructing
the documents and the programs; every commit of the repository says
so and carries the prompts that led to it, verbatim.}%
\thanks{\buildstamp}}

\markboth{A Reservation Station in TxHDL}{Filmar and Janković: A Reservation Station in TxHDL}

\begin{document}
\maketitle

\begin{abstract}
A reservation station gathers the operands of an operation that
arrive on separate channels, out of order, and fires when all of
them have. TxHDL provides one as a part, \code{Station2} to
\code{Station10}, each a lowered unit with a channel per input and
a channel out. This document says what the station does, reads its
rule cycle by cycle on the run of an example, shows how a station is
instantiated and fed and how a FIFO goes behind it, and shows the
netlist the lowering makes of it.
\end{abstract}

\section{What It Is}
\label{sec:s-what}

An operation that needs several operands waits for the last of them.
When the operands come from different producers, on different
channels and in no fixed order, something must hold the ones that
have arrived until the rest do, and must keep the operands of one
operation apart from the next one's. That something is a
reservation station, after Tomasulo. Each operation has a tag; each
operand carries the tag of its operation; the station has a line per
tag, with a cell per input; an operand lands in its cell of its line;
a line whose every cell holds is complete and goes out, all its
operands at once, with its tag.

In TxHDL the inputs are channels and so is the output. An input's
\code{ready} is whether the station can take what it offers, which
is whether the cell the offer's tag names is free. The output's
\code{valid} is a complete line, and its \code{ready} is the room
the channel behind it has.

\section{The Interface}
\label{sec:s-interface}

An input carries a \code{Tagged<TB, T>}, the tag of \code{TB} bits
above a value of type \code{T}, a struct that derives
\code{Transaction} and \code{Value} as any channel payload does. A
station of \code{N} inputs is \code{StationN<TB, L, T0, .., TN>}:
the tag width, the line count, which is \code{1 << TB} and is
stated since a computed width in a type needs nightly Rust, and one
value type per input. Its \code{run} takes a tuple of \code{N}
receivers, \code{Rx<Tagged<TB, Ti>>}, and one sender,
\code{Tx<LineN<TB, T0, .., TN>>}, where \code{LineN} is the tag and
the values, \code{v0} to \code{vN}, a derived struct beside the
station. Listing~\ref{lst:s-part} is the file that provides them:
the payload struct, and the nine stations written out by
\code{station!}, one per count of inputs, because the lowering reads
a body and not a loop over inputs.

\lstinputlisting[rangebeginprefix=//\ begin\{,rangebeginsuffix=\},rangeendprefix=//\ end\{,rangeendsuffix=\},includerangemarker=false,linerange=part-part,caption={From \code{lib/parts/src/station.rs}: the tagged payload and the nine stations; the rule is in the module's own words there, and its tests follow.},label={lst:s-part}]{lib/parts/src/station.rs}

\section{The Rule}
\label{sec:s-rule}

Every cycle the station looks at what each input offers, its tag and
its value, and at its own state: a memory per input holding the
cells, and a mask per input with one bit per line saying which cells
hold. Then:

\begin{enumerate}
\item An input's cell is \emph{free} when its tag's bit of the
  input's mask is clear.
\item Taking an input \emph{completes} its line when, for every
  other input, the cell of that line holds, or that input offers the
  same tag now and its own cell is free.
\item The line \emph{sent} is the completing input's of lowest
  index, if any completes and the output has room. At most one line
  leaves per cycle.
\item An input is \emph{taken} when its cell is free and taking it
  completes nothing, or completes the line being sent. An input whose
  taking would complete a line that is not sent this cycle waits,
  which is the case when the output has no room, and when another
  input of lower index completes a different line at the same time.
\item A taken input writes its cell and sets its bit. The line sent
  clears its bit in every mask, and its values come from the cells,
  or, for an input completing the line this cycle, straight from
  that input.
\end{enumerate}

So an input never waits on another input's cell, only on its own,
or on the output's room. What the rule does not do is protect
against a stream that never completes a line: if an input's next
offer is a tag whose cell it already holds, it waits until that line
completes, and if the other inputs never offer that tag, it waits
forever, with everything behind it in its channel. A station is not
an arbiter; each input must, in time, see every tag its lines wait
on.

\section{The Rule, Cycle by Cycle}
\label{sec:s-cycles}

The example is a \code{Station2} of a byte and a half word by a
two-bit tag, so four lines. Its first source offers bytes tagged
0, 1, 2, 3, 0, 2, one per cycle while the station takes; its second
offers half words tagged 2, 0, 3, 1, 2, 0, skipping every third
cycle; and the sink takes the output on two cycles in four.
Listing~\ref{lst:s-out} is what the run printed, one row per cycle:
what each source offered, the two masks as the cycle left them, and
the line the sink took. An offer is seen by the station the cycle
after it is made, since a send lands at the end of its step, and a
line sent is taken the cycle after at the earliest, for the same
reason.

\lstinputlisting[style=ascii,caption={What \code{ex\_station} printed when the build ran it: the offers, the masks, the lines taken; then the lowering.},label={lst:s-out}]{out_station.txt}

Read down the rows.

\begin{itemize}
\item Cycle 1: the byte tagged 0 and the half word tagged 2 are
  offered. Cycle 2: the station takes both, into cell 0 of the first
  mask and cell 2 of the second, which the row for cycle 2 shows as
  \code{0001} and \code{0100}.
\item Cycle 2 offered the byte tagged 1 and the half word tagged 0.
  In cycle 3 the half word's tag 0 completes line 0, since the
  byte's cell 0 holds: line 0 is sent, the half word goes straight
  through, and cell 0 of the first mask is cleared as the byte
  tagged 1 fills cell 1: \code{0010}. The sink takes line 0 in cycle
  4, the row that shows it.
\item Cycle 3 offered the byte tagged 2. In cycle 4 it completes
  line 2, since the half word's cell 2 holds; line 2 is sent, cell 2
  of the second mask is cleared: \code{0000}.
\item Cycle 4 offered a byte and a half word both tagged 3. In cycle
  5 each completes the line through the other's offer, both are
  taken, line 3 is sent at once, and neither mask changes.
\item Cycle 5 offered the byte tagged 0 and the half word tagged 1.
  In cycle 6 the byte fills cell 0, \code{0011}; the half word's tag
  1 would complete line 1, since the byte's cell 1 holds, but the
  output holds lines 2 and 3, which the sink has not taken, so there
  is no room and the half word waits. In cycle 7 it waits again; the
  byte tagged 2 offered in cycle 6 fills cell 2, \code{0111}. In
  cycle 8 the sink has taken line 2, there is room, line 1 is sent,
  and cell 1 clears: \code{0101}.
\item The half word tagged 2, offered in cycle 7, completes line 2
  in cycle 9, \code{0001}; the last half word, tagged 0, completes
  line 0 in cycle 11, and the mask is empty. The sink takes lines 1,
  2 and 0 in cycles 11, 12 and 15.
\end{itemize}

Figure~\ref{fig:s-wave} is the same run as a waveform: the inputs'
tags with their \code{valid} and \code{ready}, the masks, and the
output's tag and values with its handshake. The \code{ready} of the
second input is low through cycles 6 and 7, the wait of the rule's
fourth point.

\begin{figure*}[t]
\centering
\input{station_timing.tex}
\caption{The example's run: the two inputs' tags, values and
handshakes, the masks, and the lines out with their values.}
\label{fig:s-wave}
\end{figure*}

\section{Using a Station}
\label{sec:s-use}

Listing~\ref{lst:s-example} is the example whole. The station is a
type, \code{Station2<2, 4, U<8>, U<16>>}, named once; its inputs are
the receiving ends of two channels and its output the sending end of
a third; the sources send \code{Tagged} values and the sink receives
\code{Line2} values, reading \code{tag}, \code{v0} and \code{v1}. A
station of three or more inputs is the same with more channels,
\code{Station3<2, 4, A, B, C>} and so on to ten; the value types may
differ, the tag width is one for all. The tag has two bits at least,
since a one-bit value is a single wire in VHDL and cannot index a
line.

\lstinputlisting[caption={\code{ex\_station.rs}. Run with \code{bazel run //lib/examples:ex\_station}.},label={lst:s-example}]{lib/examples/ex_station.rs}

A channel holds two, so the station's output has room for two lines
before it waits. Where the consumer takes in bursts, or holds off
longer, the \code{Fifo} part goes behind the station: a channel in,
a channel out and \code{D} words between, \code{D} a power of two
stated with its address width, \code{Fifo<Line2<2, U<8>, U<16>>,
3, 8>} for eight lines. Listing~\ref{lst:s-queue} is the FIFO's own
example run: a burst of nine words into a FIFO of four, the fill,
the hold of the source, and the drain in order.

\lstinputlisting[style=ascii,caption={What \code{ex\_queue} printed when the build ran it: a FIFO of four between two channels, filling and draining; then its lowering.},label={lst:s-queue}]{out_queue.txt}

\section{The Same Station in Verilog}
\label{sec:s-verilog}

Listing~\ref{lst:s-v} is the station of the example written by hand
in Verilog, for a reader who knows what that looks like: two inputs,
a byte and a half word by a two-bit tag, four lines, the same ports
the lowering gives the TxHDL one. It is not a translation of the
netlist; it is what an engineer writes from the rule of
Section~\ref{sec:s-rule}, and it runs beside the TxHDL station in the
example, as a foreign unit fed the same offers, with the build
asserting that every line the two send agrees; Listing~\ref{lst:s-out}
ends with the count. Listing~\ref{lst:s-src} is the other side: the
text \code{station!} wrote for the two-input station, printed by the
example, since the macro keeps its text. Listing~\ref{lst:s-sizes}
counts the lines of both.

\lstinputlisting[style=ascii,caption={\code{station.v}: the two-input station by hand.},label={lst:s-v}]{lib/examples/station.v}

\lstinputlisting[caption={\code{station2.rs}: the two-input station as \code{station!} wrote it, printed by \code{ex\_station --source}.},label={lst:s-src}]{station2.rs}

\lstinputlisting[style=ascii,caption={The two, by lines, as the build counted them.},label={lst:s-sizes}]{station_sizes.txt}

Read side by side, the two say the same hardware, and the wires have
the same names, since the rule has words for them. What differs is
what each author had to write, and what each got checked.

\begin{itemize}
\item \textbf{Widths and slices.} The Verilog names every width,
  \code{[9:0]}, \code{[17:0]}, \code{[25:0]}, and slices the tag and
  the value out of each input by hand, \code{in0\_data[9:8]}; a
  change of a value type is a change in five places. The TxHDL
  station takes them from the payload type: \code{head0.tag} is the
  field, wherever the derive laid it, and the value types are
  parameters of the type. The Verilog could take widths as
  parameters too; what it cannot take as a parameter is the count of
  inputs, which is why the Verilog is one station and the macro
  writes nine.
\item \textbf{The handshake.} Each \code{ready} and \code{valid}
  is a wire the Verilog author writes and gets right, or not: the
  take must be \code{valid} and \code{ready} together, with no
  combinational path from an output back to an input's decision that
  a placer would refuse. In the TxHDL station \code{recv\_if(take0)}
  is the take, and the lowering makes the \code{ready} of it; a
  \code{send} under \code{if} is the \code{valid}. The author writes
  what is taken and what is sent, not the wires.
\item \textbf{The check.} The Verilog has no test until someone
  writes a testbench, and this one has one only because it runs
  beside the TxHDL station. The TxHDL station ran as a program
  against the rule written with loops for four hundred cycles before
  it was ever a netlist, and its netlist is checked against its own
  run under two simulators on every build, in both languages.
\item \textbf{Lines.} By the count in Listing~\ref{lst:s-sizes} the
  Verilog is the shorter text, by about half once comments are left
  out. The TxHDL text is a generator's, one clause per line, and it
  carries what the Verilog does not: the line struct, the types as
  parameters, and the struct of the state that the trace and the
  netlist read. Count instead what an author writes: the Verilog
  author writes every line of the fifty, for this count of inputs
  and these widths; the TxHDL designer writes \code{station!(Station2,
  2)} once for all widths, and the type at the use.
\end{itemize}

\section{The Netlist}
\label{sec:s-netlist}

The end of Listing~\ref{lst:s-out} is the Verilog of the example's
station, as the lowering wrote it. Each mask is a register of four
bits and each input's cells a memory of four words; every wire of
the step is named for the rule's word: \code{offered0},
\code{cell\_free0}, \code{completes0}, \code{line\_tag},
\code{send\_line}, \code{take0}, \code{cell\_bit0},
\code{line\_clear}, \code{out0}. The tag and the value of an input
are slices of the input's data, at the bits the payload's derive laid
them out; the one-hot of a cell is a one shifted by the tag; the
masks are updated in one assignment each, the cell filled or'd in
and the line sent and'd out; the output's data is the concatenation
of the tag and the values, its \code{valid} the \code{send\_line}
wire, and each input's \code{ready} its \code{take}. The VHDL says
the same with \code{shift\_left} and \code{to\_integer}. Both are
simulated against the run above, by the build, every time.

\begin{thebibliography}{9}

\bibitem{tutorial}
F.~Filmar and D.~Janković, ``TxHDL, step by step,'' project
repository \code{HDL/txhdl}, \code{//docs:tutorial}, 2026.

\bibitem{runtime}
F.~Filmar and D.~Janković, ``The TxHDL runtime,'' project repository
\code{HDL/txhdl}, \code{//docs:runtime}, 2026.

\bibitem{examples}
F.~Filmar and D.~Janković, ``TxHDL by example,'' project repository
\code{HDL/txhdl}, \code{//docs:examples}, 2026.

\bibitem{cover}
F.~Filmar and D.~Janković, ``TxHDL documents,'' project repository
\code{HDL/txhdl}, \code{//docs:cover}, 2026.

\end{thebibliography}

\end{document}
